% Part: sets-functions-relations
% Chapter: arithmetization
% Section: cuts
%
\documentclass[../../../include/open-logic-section]{subfiles}

\begin{document}

\olfileid{sfr}{arith}{cuts}
\olsection{$\Rat$ पासून $\Real$ कडे}

मूलतः, संपूर्णता गुणधर्म असे दाखवतो की वास्तव संख्या-रेषेवरील कोणताही
$\alpha$ बिंदू त्या रेषेचे अचूक दोन भाग करतो: खालच्या भागाचा $\alpha$
हा लघुतम ऊर्ध्वबंध असतो आणि वरच्या भागाचा $\alpha$ हा महत्तम निम्नबंध
असतो. परिमेय संख्यांपासून वास्तव संख्यांची \emph{रचना} करण्यासाठी
डेडेकिंड यांनी वास्तव संख्यांना परिमेय संख्यांचे विभाजन करणारे
\emph{छेद} मानण्याची सूचना केली. म्हणजेच, $< \sqrt{2}$ असलेल्या परिमेय
संख्यांना $> \sqrt{2}$ असलेल्या परिमेय संख्यांपासून वेगळे करणाऱ्या
\emph{छेदाशी} आपण $\sqrt{2}$ ची एकरूपता मानतो.

ही मांडणी नीट नेमकी करू या. परिमेय संख्यांचे दोन भाग केले, तर आपण केलेले
विभाजन केवळ त्याचा \emph{खालचा} भाग विचारात घेऊन अनन्यपणे ओळखता येते.
म्हणून नेमकी पुढील व्याख्या देऊ:

\begin{defn}[छेद]
\emph{छेद} $\alpha$ म्हणजे परिमेय संख्यांचा असा कोणताही अरिक्त उचित
आरंभीचा खंड ज्याला महत्तम घटक नाही. म्हणजे $\alpha$ हा छेद तेव्हा आणि
केवळ तेव्हाच असतो, जेव्हा:
\begin{enumerate}
	\item \emph{अरिक्त, उचित}: $\emptyset \neq \alpha \subsetneq \Rat$
	\item \emph{आरंभीचा}: सर्व $p,q \in \Rat$ साठी: $p < q \in \alpha$ असेल, तर $p \in \alpha$
	\item \emph{महत्तम घटक नाही}: प्रत्येक $p \in \alpha$ साठी $p < q$ असे काही $q \in \alpha$ असते
\end{enumerate}
मग $\Real$ हा छेदांचा संच आहे.
\end{defn}

आता $\sqrt{2} = \Setabs{p \in \Rat}{p^2 < 2\text{
किंवा }p < 0}$ असे म्हणता येते. अर्थात, हा खरोखर \emph{छेद} आहे हे
तपासले पाहिजे; ती तपासणी आपण \olref[check]{sec} कडे सोपवतो.

आधीप्रमाणेच, काही वस्तूंची व्याख्या केल्यानंतर त्यांच्यावरील मूलभूत
फलने आणि संबंध व्याख्यायित करावे लागतात. एका सोप्या व्याख्येपासून सुरुवात करू:
\begin{align*}
	\alpha \leq \beta \text{ तेव्हा आणि केवळ तेव्हाच }\alpha \subseteq \beta
\end{align*}
क्रमाची ही व्याख्या छेदांच्या संचाला संपूर्णता गुणधर्म आहे हा मध्यवर्ती
निष्कर्ष \emph{मांडू} देते. पूर्णपणे उलगडल्यास विधानाचे रूप असे आहे.
$S$ हा ऊर्ध्वबंध असलेला छेदांचा अरिक्त संच असेल, तर $S$ ला लघुतम
ऊर्ध्वबंध असतो. अधिक तपशीलाने: $S$ चा ऊर्ध्वबंध असणारा $\lambda$ हा छेद
असतो, म्हणजे $(\forall \alpha \in S)\alpha \subseteq \lambda$, आणि
$\lambda$ हा अशा छेदांपैकी लघुतम असतो, म्हणजे $(\forall \beta \in \Real)((\forall \alpha \in S)\alpha \subseteq \beta \lif \lambda \subseteq \beta)$. आता या निष्कर्षाची सिद्धता पाहा:

\begin{thm}\ollabel{realcompleteness}
छेदांच्या संचाला संपूर्णता गुणधर्म आहे.
\end{thm}

\begin{proof}
$S$ हा ऊर्ध्वबंध असलेला छेदांचा कोणताही अरिक्त संच असू द्या. $\lambda
= \bigcup S$ असे ठरवू.
%\Setabs{p \in \Rat}{(\exists \alpha \in S)p \in \alpha}$$
प्रथम $\lambda$ हा छेद आहे असा दावा करू:
\begin{enumerate}
\item $S$ अरिक्त असल्यामुळे $S$ मध्ये किमान एक छेद आहे, म्हणून
$\emptyset \neq \lambda$. $S$ हा छेदांचा संच असल्यामुळे $\lambda
\subseteq \Rat$. $S$ ला ऊर्ध्वबंध असल्यामुळे प्रत्येक $\alpha \in S$
मधून अनुपस्थित असलेली अशी काही $p \in \Rat$ आहे. म्हणून $p\notin \lambda$,
आणि परिणामी $\lambda \subsetneq \Rat$.
%READERNOTE{MRARITH-004: गोठवलेल्या इंग्रजी सिद्धतेत λ अरिक्त असल्याचे पहिले कारण चुकून “S has an upper bound” असे दिले आहे. आवश्यक आणि आधीच गृहीत असलेले कारण S अरिक्त असणे हे आहे; मराठी सिद्धतेत तेच वापरले आहे. पुढील उचित-उपसंच युक्तिवादासाठी मात्र ऊर्ध्वबंधाचे गृहीतक योग्य आहे.}
\item $p < q \in \lambda$ असे समजा. मग $q \in \alpha$ असे काही
$\alpha \in S$ आहे. $\alpha$ हा छेद असल्यामुळे $p \in \alpha$.
म्हणून $p \in \lambda$.
\item $p \in \lambda$ असे समजा. मग $p \in \alpha$ असे काही
$\alpha \in S$ आहे. $\alpha$ हा छेद असल्यामुळे $p < q$ असे काही
$q \in \alpha$ आहे. म्हणून $q \in \lambda$.
\end{enumerate}
यावरून दावा सिद्ध होतो. शिवाय, स्पष्टपणे $(\forall \alpha \in S)\alpha
\subseteq \bigcup S = \lambda$, म्हणजे $\lambda$ हा $S$ चा ऊर्ध्वबंध
आहे. आता $\beta \in \mathbb{R}$ हाही ऊर्ध्वबंध आहे असे समजा, म्हणजे
$(\forall \alpha \in S)\alpha \subseteq \beta$. कोणत्याही $p \in \Rat$
साठी, $p \in \lambda$ असेल तर $p \in \alpha$ असे काही $\alpha \in S$
आहे, आणि त्यामुळे $p \in \beta$. सार्वत्रिकीकरण केल्यास $\lambda
\subseteq \beta$. म्हणून $\lambda$ हा $S$ चा \emph{लघुतम} ऊर्ध्वबंध आहे.
\end{proof}

म्हणून आपल्याकडे संपूर्णता गुणधर्म पूर्ण करणाऱ्या अनेक वस्तू आहेत.
हे सांगण्याचा एक मार्ग असा: आपल्या छेदांमध्ये ``पोकळ्या'' नाहीत.
(त्यामुळे परिमेय संख्यांऐवजी वास्तव संख्यांचे आणखी ``छेद'' घेतल्यास
काही रोचक नवी वस्तू मिळणार नाही.)

पुढे वास्तव संख्यांवरील काही क्रिया व्याख्यायित केल्या पाहिजेत. प्रत्येक
$p \in \Rat$ साठी $p_\Real = \Setabs{q \in \Rat}{q < p}$ असे ठरवून
परिमेय संख्यांना वास्तव संख्यांत अंतःस्थापित करण्यापासून सुरुवात करू.
मग पुढील व्याख्या करतो:
\begin{align*}
	\alpha + \beta &= \Setabs{p + q}{p \in \alpha \land q \in \beta}\\
%	\alpha - \beta &\defis \Setabs{p - q}{p \in \alpha \land q \in \Rat \setminus \beta}\\
	\alpha \times \beta &=
	\Setabs{p \times q}{0 \leq p \in \alpha \land 0 \leq q \in \beta} \cup 0_\Real & \text{जर }\alpha, \beta \geq 0_\Real
%	\alpha \div \beta &\defis
%	\Setabs{\nicefrac{p}{q}}{p \in \alpha \land q \in \Rat \setminus \beta}  & \text{जर }\alpha \geq 0_\Real, \beta > 0_\Real
\end{align*}
%READERNOTE{MRARITH-005: गोठवलेल्या इंग्रजी गुणाकार-व्याख्येत शून्याचा छेद एकदाच 0^R असा विसंगत ऊर्ध्वनिर्देशांक वापरून लिहिला आहे; त्याच परिच्छेदातील सर्व नियंत्रक संकेतन 0_R वापरते. मराठी सूत्रात 0_R वापरले आहे.}
गुणाकाराची इतर प्रकरणे हाताळण्यासाठी प्रथम $0_\Real \times \alpha = 0_\Real = \alpha \times 0_\Real$ असे ठरवू आणि पुढील व्याख्या करू:
%READERNOTE{MRARITH-006: गोठवलेल्या इंग्रजी स्रोताने शून्यासोबतच्या गुणाकाराची ही आवश्यक व्याख्या त्याच ओळीवर निष्क्रिय केली आहे; त्यामुळे एक घटक ऋण व दुसरा शून्य असलेली प्रकरणे पुढील प्रकरणनिहाय व्याख्येत येत नाहीत. स्रोताच्याच निष्क्रिय सूत्राला मराठीत सक्रिय केले आहे.}
\begin{align*}
	-\alpha &= \Setabs{p - q}{p < 0 \land q \notin \alpha}
\end{align*}
आणि मग पुढीलप्रमाणे ठरवू:
\begin{align*}
	\alpha \times \beta &\defis
	\begin{cases}
		\mathord{-}\alpha \times \mathord{-}\beta &\text{जर }\alpha < 0_\Real\text{ आणि }\beta < 0_\Real\\
		\mathord{-}(\mathord{-}\alpha \times \beta) &\text{जर }\alpha < 0_\Real \text{ आणि }\beta > 0_\Real\\
		\mathord{-}(\alpha \times \mathord{-}\beta) &\text{जर }\alpha > 0_\Real \text{ आणि }\beta < 0_\Real
	\end{cases}
\end{align*}
यानंतर यांपैकी प्रत्येक व्याख्येतून नेहमी छेदच मिळतो हे तपासावे लागेल.
आणि शेवटी, अशा प्रकारे व्याख्यायित केलेले छेद अपेक्षेप्रमाणेच वागतात
हे दाखवणारी सोपी (पण लांबलचक) सिद्धता पूर्ण केली पाहिजे. मात्र हे सर्व
आपण \olref[check]{sec} कडे सोपवतो.
\end{document}
